% This document was generated by an LLM.

\documentclass{essay}

\title{Ordinary Differential Equations}
\subtitle{A Reference List of Types and Techniques}
\author{}
\date{}

\begin{document}

\maketitle

\section*{How to use this document}

The central skill in solving ODEs is \emph{recognition}: before reaching for a method, identify the structural type of the equation.  Each section below names a type, states the recognition cue, and gives worked examples.  The sections are ordered so that each technique builds naturally on the previous ones.

The document is divided into two parts.  Part~I covers closed-form methods: you recognise a pattern and obtain an elementary function as the answer.  Part~II covers series methods, which you reach when Part~I offers no foothold.

\part{Closed-Form Methods}

\section{Separable equations}

\textbf{Recognition.}  The equation can be written as
\[
  f(y)\,dy = g(x)\,dx,
\]
i.e.\ the right-hand side factors into a function of $x$ alone times a function of $y$ alone.

\textbf{Method.}  Separate and integrate both sides independently.

\textbf{Examples.}
\[
  \frac{dy}{dx} = xy^2, \qquad
  \frac{dy}{dx} = \frac{1+y^2}{1+x^2}, \qquad
  \frac{dy}{dx} = y\cos x.
\]

\section{First-order linear equations (integrating factor)}

\textbf{Recognition.}  The equation is linear in $y$ and $y'$:
\[
  y' + P(x)\,y = Q(x).
\]

\textbf{Method.}  Multiply through by the integrating factor $\mu(x) = e^{\int P(x)\,dx}$.  The left side becomes $(\mu y)'$; integrate.

\textbf{Examples.}
\[
  y' + 2y = e^{-x}, \qquad
  y' + \frac{y}{x} = x^2, \qquad
  y' - y\tan x = \sin x.
\]

\section{Exact equations}

\textbf{Recognition.}  Written as $M(x,y)\,dx + N(x,y)\,dy = 0$, and
\[
  \frac{\partial M}{\partial y} = \frac{\partial N}{\partial x}.
\]

\textbf{Method.}  Find $F$ satisfying $F_x = M$ and $F_y = N$.  The solution is $F(x,y) = C$.

\textbf{Examples.}
\[
  (2xy + 3)\,dx + (x^2 - 1)\,dy = 0, \qquad
  e^y\,dx + (xe^y + 2y)\,dy = 0.
\]

\section{Non-exact equations made exact (integrating factor for forms)}

\textbf{Recognition.}  The equation $M\,dx + N\,dy = 0$ is not exact, but
\[
  \frac{M_y - N_x}{N} \text{ depends only on } x
  \quad \text{(or } \frac{N_x - M_y}{M} \text{ on } y \text{)}.
\]

\textbf{Method.}  Compute $\mu$ from the ratio above and multiply through to make the equation exact.

\textbf{Example.}
\[
  (2y^2 + 3x)\,dx + 2xy\,dy = 0.
\]

\section{Homogeneous equations (substitution $v = y/x$)}

\textbf{Recognition.}  Both sides are homogeneous of the same degree, or equivalently
$dy/dx = F(y/x)$ for some function $F$.

\textbf{Method.}  Substitute $y = vx$, so $y' = v + xv'$.  The equation becomes separable in $v$ and $x$.

\textbf{Examples.}
\[
  \frac{dy}{dx} = \frac{x+y}{x}, \qquad
  (x^2 + y^2)\,dx - 2xy\,dy = 0.
\]

\section{Bernoulli equations}

\textbf{Recognition.}  The equation has the form
\[
  y' + P(x)\,y = Q(x)\,y^n, \quad n \neq 0, 1.
\]
It is nonlinear because of the $y^n$ term on the right.

\textbf{Method.}  Substitute $v = y^{1-n}$ to obtain a linear equation in $v$.

\textbf{Examples.}
\[
  y' + \frac{y}{x} = xy^2, \qquad y' + y = y^3.
\]

\section{Substitution $v = ax + by + c$}

\textbf{Recognition.}  The right-hand side depends only on the linear combination $ax+by+c$:
\[
  \frac{dy}{dx} = f(ax + by + c).
\]

\textbf{Method.}  Set $v = ax + by + c$.  Then $v' = a + b f(v)$, which is separable.

\textbf{Examples.}
\[
  \frac{dy}{dx} = (x + y + 1)^2, \qquad \frac{dy}{dx} = \cos(x+y).
\]

\section{Second-order linear, constant coefficients, homogeneous}

\textbf{Recognition.}  The equation is
\[
  ay'' + by' + cy = 0, \quad a,b,c \in \R, \; a \neq 0.
\]

\textbf{Method.}  Solve the characteristic equation $ar^2 + br + c = 0$ and apply the appropriate case.

\subsection*{Case 1: distinct real roots $r_1 \neq r_2$}
\[
  y'' - 5y' + 6y = 0 \quad (r = 2, 3)
  \;\Rightarrow\; y = C_1 e^{2x} + C_2 e^{3x}.
\]

\subsection*{Case 2: repeated root $r_1 = r_2 = r$}
\[
  y'' - 4y' + 4y = 0 \quad (r = 2)
  \;\Rightarrow\; y = (C_1 + C_2 x)e^{2x}.
\]

\subsection*{Case 3: complex conjugate roots $r = \alpha \pm \beta i$}
\[
  y'' + 2y' + 5y = 0 \quad (r = -1 \pm 2i)
  \;\Rightarrow\; y = e^{-x}(C_1\cos 2x + C_2 \sin 2x).
\]

\section{Method of undetermined coefficients}

\textbf{Recognition.}  The equation is $ay'' + by' + cy = g(x)$ where $g$ is a polynomial, exponential, sine/cosine, or a product of these.

\textbf{Method.}  Guess a particular solution $y_p$ of the same form as $g$.  If any term of the trial $y_p$ solves the homogeneous equation, multiply the trial form by $x$ (or $x^2$, etc.) to avoid resonance.

\textbf{Examples.}
\[
  y'' - 3y' + 2y = e^{3x}, \qquad
  y'' + y = \sin 2x, \qquad
  y'' - y = x^2.
\]
\textbf{Resonance example:} $y'' + y = \cos x$ — the trial form $A\cos x$ solves the homogeneous equation, so use $y_p = x(A\cos x + B\sin x)$ instead.

\section{Variation of parameters}

\textbf{Recognition.}  The equation is $y'' + P(x)y' + Q(x)y = g(x)$, and undetermined coefficients is not applicable because $g$ is not of the required form (e.g.\ $g = \sec x$, $\tan x$, $\ln x$).

\textbf{Method.}  Let $y_1, y_2$ be a fundamental pair of solutions to the homogeneous equation with Wronskian $W$.  Then
\[
  y_p = -y_1 \int \frac{y_2\, g}{W}\,dx + y_2 \int \frac{y_1\, g}{W}\,dx.
\]

\textbf{Examples.}
\[
  y'' + y = \sec x, \qquad y'' - 2y' + y = \frac{e^x}{x}.
\]

\section{Cauchy--Euler (equidimensional) equations}

\textbf{Recognition.}  The coefficient of $y^{(k)}$ is a constant multiple of $x^k$:
\[
  x^2 y'' + \alpha x y' + \beta y = g(x).
\]

\textbf{Method.}  Try $y = x^r$.  This yields a quadratic indicial equation in $r$, with the same three-case structure as the constant-coefficient characteristic equation.  Alternatively, substitute $x = e^t$ to convert to a constant-coefficient equation in $t$.

\textbf{Examples.}
\[
  x^2 y'' + 2xy' - 6y = 0, \qquad x^2 y'' - 3xy' + 4y = x^2.
\]

\section{Reduction of order}

\textbf{Recognition.}  A second-order linear ODE for which one solution $y_1$ is already known.

\textbf{Method.}  Set $y = v(x)\,y_1(x)$.  Substituting reduces the order by one (the $v$ term drops out), yielding a first-order linear equation in $v'$.

\textbf{Example.}
\[
  x^2 y'' - xy' + y = 0, \quad y_1 = x.
\]

\section{Missing dependent variable ($y'' = f(x,\, y')$)}

\textbf{Recognition.}  The equation does not contain $y$ explicitly.

\textbf{Method.}  Set $p = y'$.  Then $y'' = p'$, reducing to a first-order equation in $p(x)$.

\textbf{Example.}
\[
  y'' = x + y'.
\]

\section{Missing independent variable ($y'' = f(y,\, y')$)}

\textbf{Recognition.}  The equation does not contain $x$ explicitly.

\textbf{Method.}  Set $p = y'$ and write
\[
  y'' = \frac{dp}{dx} = \frac{dp}{dy}\frac{dy}{dx} = p\,\frac{dp}{dy}.
\]
This yields a first-order equation in $p(y)$.

\textbf{Examples.}
\[
  y\,y'' + (y')^2 = 0, \qquad y'' + (y')^2 = 0.
\]

\part{Series Methods}

These methods apply when Part~I offers no foothold: variable coefficients that are not equidimensional, no known closed-form solution, no obvious second solution for reduction of order.

\medskip
\noindent\textbf{Structural parallel.}  The indicial equation in the Frobenius method plays the same role as the characteristic equation in constant-coefficient ODEs.  Its two roots $r_1, r_2$ determine solution form, and the three cases below mirror exactly the three root cases of Section~8.

\begin{center}
  \renewcommand{\arraystretch}{1.4}
  \begin{tabular}{ll}
    \toprule
    \textbf{Constant-coefficient char.\ eq.} & \textbf{Frobenius indicial eq.} \\
    \midrule
    Distinct real roots & $r_1 - r_2 \notin \mathbb{Z}$: two clean Frobenius series \\
    Repeated root & $r_1 = r_2$: one series, second solution has $\ln x$ \\
    Complex conjugate roots & $r_1 - r_2 \in \mathbb{Z}_{>0}$: second solution \emph{may} have $\ln x$ \\
    \bottomrule
  \end{tabular}
\end{center}

The repeated-root logarithm arises in Frobenius for the same reason a factor of $x$ appears in the constant-coefficient case: a dimension of solution space must be manufactured by differentiating with respect to the parameter $r$.

\section{Power series at an ordinary point}

\textbf{Recognition.}  Write the equation as $y'' + p(x)y' + q(x)y = 0$.  If $p$ and $q$ are analytic at $x_0$ (e.g.\ polynomial coefficients with leading coefficient nonzero at $x_0$), then $x_0$ is an \emph{ordinary point}.  Two linearly independent power series solutions converge in a neighbourhood.

\textbf{Method.}  Substitute $y = \sum_{n=0}^{\infty} a_n (x - x_0)^n$, collect powers of $(x-x_0)$, and derive a recurrence relation for $a_n$.

\subsection*{Warm-ups (recovery of known functions)}

These confirm the bookkeeping before meeting genuinely new functions.
\[
  y' = y \;\Rightarrow\; y = \sum_{n=0}^{\infty} \frac{x^n}{n!} = e^x.
\]
\[
  y'' + y = 0 \;\Rightarrow\; \text{two series recovering } \cos x \text{ and } \sin x.
\]

\subsection*{Airy equation}

\[
  y'' - xy = 0.
\]

The simplest equation whose solutions are genuinely non-elementary.  The recurrence links every third coefficient; the two independent solutions are the Airy functions $\mathrm{Ai}(x)$ and $\mathrm{Bi}(x)$.  Work this one by hand at least once.

\subsection*{Hermite equation}

\[
  y'' - 2xy' + 2\lambda y = 0.
\]

When $\lambda = n$ is a non-negative integer, one of the two series terminates, producing the Hermite polynomial $H_n(x)$.  The ``termination at integer parameter'' pattern recurs throughout classical ODEs.

\subsection*{Legendre equation}

\[
  (1-x^2)y'' - 2xy' + n(n+1)y = 0.
\]

Solved around the ordinary point $x_0 = 0$.  One series terminates when $n$ is a non-negative integer, giving the Legendre polynomials $P_n(x)$.  Note that $x = \pm 1$ are singular points, making them natural candidates for Part~B below.

\subsection*{Chebyshev equation}

\[
  (1-x^2)y'' - xy' + \lambda^2 y = 0.
\]

Same singular-point structure as Legendre.  Polynomial solutions arise when $\lambda = n$.

\section{Frobenius method at a regular singular point}

\textbf{Recognition.}  The point $x_0$ is singular (the leading coefficient vanishes there), but
\[
  (x - x_0)\,p(x) \quad\text{and}\quad (x-x_0)^2\,q(x)
\]
are both analytic at $x_0$.  Such a point is called a \emph{regular} singular point.

\textbf{Method.}  Substitute $y = (x-x_0)^r \sum_{n=0}^{\infty} a_n (x-x_0)^n$ with $a_0 \neq 0$.  The lowest power of $(x-x_0)$ yields the \textbf{indicial equation}
\[
  r(r-1) + p_0\,r + q_0 = 0,
\]
where $p_0 = \lim_{x\to x_0}(x-x_0)p(x)$ and $q_0 = \lim_{x\to x_0}(x-x_0)^2 q(x)$.  Denote the roots $r_1 \geq r_2$.

\subsection*{Case 1: $r_1 - r_2 \notin \mathbb{Z}$}

Both roots give independent Frobenius series with no logarithms.

\textbf{Example.}
\[
  2x\,y'' + y' + y = 0.
\]
Indicial equation: $r(2r-1) = 0$, roots $r_1 = \tfrac{1}{2}$, $r_2 = 0$.  Difference $\tfrac{1}{2}$, not an integer.  Two clean series follow.

\subsection*{Case 2: $r_1 = r_2$ (repeated root)}

One Frobenius series comes from $r_1$.  The second solution necessarily involves a logarithm:
\[
  y_2 = y_1 \ln x + x^{r_1}\sum_{n=1}^{\infty} b_n x^n.
\]

\textbf{Example: Bessel equation of order zero.}
\[
  x^2 y'' + xy' + x^2 y = 0.
\]
Indicial equation: $r^2 = 0$, so $r_1 = r_2 = 0$.  The first solution is $J_0(x)$.  The second solution $Y_0(x)$ contains $J_0(x)\ln x$ plus a Frobenius-type series.

\subsection*{Case 3: $r_1 - r_2 \in \mathbb{Z}_{>0}$ (integer-differing roots)}

The larger root $r_1$ always produces a Frobenius series.  The smaller root $r_2$ \emph{may or may not} work: if the recurrence breaks down, the second solution requires a logarithm; if it does not break down, two clean series exist.

\medskip
\noindent\textbf{Example A (no logarithm):}
\[
  xy'' + 2y' + xy = 0,
\]
indicial roots $r = 0, -1$ (differ by 1).  The smaller root happens to work; two series are obtained without a logarithm.

\medskip
\noindent\textbf{Example B (logarithm required): Bessel equation of order one.}
\[
  x^2 y'' + xy' + (x^2 - 1)y = 0,
\]
indicial roots $r = \pm 1$ (differ by 2).  The smaller root fails in the recurrence; the second solution $Y_1(x)$ requires a logarithmic term.

\medskip
Comparing Examples A and B side-by-side is the best way to internalise why Case~3 is the subtle one.

\section{Suggested practice order}

\subsection*{Part I (closed-form methods)}

Separable $\to$ linear (integrating factor) $\to$ exact $\to$ homogeneous substitution $\to$ Bernoulli $\to$ second-order constant coefficient (all three root cases) $\to$ undetermined coefficients (include a resonance example) $\to$ Cauchy--Euler $\to$ variation of parameters $\to$ reduction of order $\to$ missing-variable substitutions.

\subsection*{Part II (series methods)}

Sanity-check series ($e^x$, then $\sin x$ and $\cos x$) $\to$ Airy equation (first genuinely new function) $\to$ Frobenius Case~1 ($2xy'' + y' + y = 0$) $\to$ Frobenius Case~2 (Bessel order 0) $\to$ Frobenius Case~3 comparison (Bessel order 0 vs.\ order 1) $\to$ Legendre and Hermite (polynomial termination at integer parameters).

\end{document}
